\documentclass[11pt,a4paper,oneside,openany]{stacks-project-book}

% Non-rendering source-note environments retained from the Stacks sources.
\usepackage{verbatim}
\newenvironment{reference}{\comment}{\endcomment}
\newenvironment{slogan}{\comment}{\endcomment}
\newenvironment{history}{\comment}{\endcomment}

% Diagrams and chapter-support packages.
\usepackage[all]{xy}
\xyoption{2cell}
\UseAllTwocells
\usepackage{multicol}
\usepackage{graphicx}

% Mainland Simplified-Chinese scientific typography.
\usepackage{fontspec}
\usepackage{xeCJK}
\defaultfontfeatures{Ligatures=TeX}
\setmainfont{Latin Modern Roman}
\setsansfont{Latin Modern Sans}
\setmonofont{Latin Modern Mono}
\setCJKmainfont[AutoFakeSlant=0.2]{Noto Serif CJK SC}
\setCJKsansfont{Noto Sans CJK SC}
\setCJKmonofont{Noto Sans CJK SC}
\newfontfamily\stackszhgreek{Noto Serif CJK SC}
\DeclareRobustCommand{\zhdelta}{{\stackszhgreek δ}}
\xeCJKsetup{PunctStyle=kaiming}
\XeTeXlinebreaklocale "zh"
\usepackage[a4paper,left=22mm,right=22mm,top=24mm,bottom=26mm,headheight=14pt,headsep=7mm,footskip=14mm]{geometry}
\linespread{1.18}
\setlength{\parindent}{2em}
\setlength{\parskip}{0pt}
\setlength{\emergencystretch}{1em}

% Localized structural names.
\renewcommand{\contentsname}{目录}
\renewcommand{\bibname}{参考文献}
\renewcommand{\proofname}{证明}
\renewcommand{\figurename}{图}
\renewcommand{\tablename}{表}
\renewcommand{\appendixname}{附录}

% Chinese numbered chapter heads while keeping numerical identifiers stable.
\makeatletter
\def\@makechapterhead#1{\global\topskip 7.5pc\relax
  \begingroup
  \fontsize{17}{22}\bfseries\centering
    \ifnum\c@secnumdepth>\m@ne
      \leavevmode \hskip-\leftskip
      \rlap{\vbox to\z@{\vss
          \ifx\chaptername\appendixname
            \centerline{\normalsize\mdseries 附录\enspace\thechapter}
          \else
            \centerline{\normalsize\mdseries 第\thechapter 章}
          \fi
          \vskip 3pc}}\hskip\leftskip\fi
     #1\par \endgroup
  \skip@34\p@ \advance\skip@-\normalbaselineskip
  \vskip\skip@ }
\makeatother

% Hyperlinks: current-volume labels resolve locally; absent chapters use tags.
\usepackage{hyperref}
\hypersetup{
  unicode=true,
  hidelinks,
  pdfborder={0 0 0},
  pdftitle={Stacks Project 简体中文版 - 阶段性累积版},
  pdfauthor={The Stacks Project contributors; zh-Hans-CN translation production},
  pdfsubject={Mainland Simplified Chinese cumulative edition of selected Stacks Project chapters}
}
\urlstyle{same}
% Generated from the frozen Stacks permanent-tag registry.
% Existing labels resolve inside this PDF; only absent labels fall back to tags.
\expandafter\def\csname stacksTag@etale-cohomology-definition-0-cech\endcsname{03NR}
\expandafter\def\csname stacksTag@etale-cohomology-definition-Cr\endcsname{03R9}
\expandafter\def\csname stacksTag@etale-cohomology-definition-G-module-continuous\endcsname{04JP}
\expandafter\def\csname stacksTag@etale-cohomology-definition-additive-sheaf\endcsname{03P4}
\expandafter\def\csname stacksTag@etale-cohomology-definition-algebraic-geometric-point\endcsname{03QX}
\expandafter\def\csname stacksTag@etale-cohomology-definition-big-etale-site\endcsname{03PI}
\expandafter\def\csname stacksTag@etale-cohomology-definition-brauer-equivalent\endcsname{03R3}
\expandafter\def\csname stacksTag@etale-cohomology-definition-brauer-group\endcsname{03R5}
\expandafter\def\csname stacksTag@etale-cohomology-definition-c\endcsname{095W}
\expandafter\def\csname stacksTag@etale-cohomology-definition-category-sheaves\endcsname{03NO}
\expandafter\def\csname stacksTag@etale-cohomology-definition-cd\endcsname{0F0Q}
\expandafter\def\csname stacksTag@etale-cohomology-definition-cd-f\endcsname{0F0Z}
\expandafter\def\csname stacksTag@etale-cohomology-definition-cech-complex\endcsname{03OL}
\expandafter\def\csname stacksTag@etale-cohomology-definition-constructible\endcsname{03RW}
\expandafter\def\csname stacksTag@etale-cohomology-definition-ctf\endcsname{03TQ}
\expandafter\def\csname stacksTag@etale-cohomology-definition-descent-datum\endcsname{03O7}
\expandafter\def\csname stacksTag@etale-cohomology-definition-descent-datum-modules\endcsname{03OB}
\expandafter\def\csname stacksTag@etale-cohomology-definition-direct-image-presheaf\endcsname{03PW}
\expandafter\def\csname stacksTag@etale-cohomology-definition-direct-image-sheaf\endcsname{03PY}
\expandafter\def\csname stacksTag@etale-cohomology-definition-effective-modules\endcsname{03OC}
\expandafter\def\csname stacksTag@etale-cohomology-definition-etale-covering\endcsname{03PG}
\expandafter\def\csname stacksTag@etale-cohomology-definition-etale-covering-initial\endcsname{03N5}
\expandafter\def\csname stacksTag@etale-cohomology-definition-etale-local-rings\endcsname{03PS}
\expandafter\def\csname stacksTag@etale-cohomology-definition-etale-morphism\endcsname{03PB}
\expandafter\def\csname stacksTag@etale-cohomology-definition-etale-topos\endcsname{04HQ}
\expandafter\def\csname stacksTag@etale-cohomology-definition-extension-zero\endcsname{03S3}
\expandafter\def\csname stacksTag@etale-cohomology-definition-family-morphisms-fixed-target\endcsname{03NG}
\expandafter\def\csname stacksTag@etale-cohomology-definition-finite-locally-constant\endcsname{03RU}
\expandafter\def\csname stacksTag@etale-cohomology-definition-fpqc-covering\endcsname{03NW}
\expandafter\def\csname stacksTag@etale-cohomology-definition-free-abelian-presheaf\endcsname{03OO}
\expandafter\def\csname stacksTag@etale-cohomology-definition-galois-cohomology\endcsname{04JR}
\expandafter\def\csname stacksTag@etale-cohomology-definition-geometric-point\endcsname{03PO}
\expandafter\def\csname stacksTag@etale-cohomology-definition-henselian\endcsname{03QF}
\expandafter\def\csname stacksTag@etale-cohomology-definition-higher-direct-images\endcsname{04I2}
\expandafter\def\csname stacksTag@etale-cohomology-definition-inverse-image\endcsname{03Q0}
\expandafter\def\csname stacksTag@etale-cohomology-definition-inverse-system-sheaves\endcsname{0EZL}
\expandafter\def\csname stacksTag@etale-cohomology-definition-locally-acyclic\endcsname{0GJP}
\expandafter\def\csname stacksTag@etale-cohomology-definition-presheaf\endcsname{03NC}
\expandafter\def\csname stacksTag@etale-cohomology-definition-ringed-site\endcsname{03OH}
\expandafter\def\csname stacksTag@etale-cohomology-definition-sheaf\endcsname{03NK}
\expandafter\def\csname stacksTag@etale-cohomology-definition-sheaf-property-fpqc\endcsname{03X6}
\expandafter\def\csname stacksTag@etale-cohomology-definition-site\endcsname{03NH}
\expandafter\def\csname stacksTag@etale-cohomology-definition-stalk\endcsname{040R}
\expandafter\def\csname stacksTag@etale-cohomology-definition-standard-etale\endcsname{03PD}
\expandafter\def\csname stacksTag@etale-cohomology-definition-standard-tau\endcsname{03X9}
\expandafter\def\csname stacksTag@etale-cohomology-definition-strictly-henselian\endcsname{03QK}
\expandafter\def\csname stacksTag@etale-cohomology-definition-structure-sheaf\endcsname{04HS}
\expandafter\def\csname stacksTag@etale-cohomology-definition-support\endcsname{04FS}
\expandafter\def\csname stacksTag@etale-cohomology-definition-tau-covering\endcsname{03X8}
\expandafter\def\csname stacksTag@etale-cohomology-definition-tau-site\endcsname{03XB}
\expandafter\def\csname stacksTag@etale-cohomology-definition-trace-map\endcsname{03SE}
\expandafter\def\csname stacksTag@etale-cohomology-definition-variety\endcsname{03RE}
\expandafter\def\csname stacksTag@etale-cohomology-equation-j-p-shriek\endcsname{0F4K}
\expandafter\def\csname stacksTag@etale-cohomology-lemma-F-1\endcsname{0A3L}
\expandafter\def\csname stacksTag@etale-cohomology-lemma-V-C-all-n-etale-fppf\endcsname{0DDS}
\expandafter\def\csname stacksTag@etale-cohomology-lemma-V-C-all-n-etale-ph\endcsname{0DE4}
\expandafter\def\csname stacksTag@etale-cohomology-lemma-affine-only-closed-points\endcsname{09AY}
\expandafter\def\csname stacksTag@etale-cohomology-lemma-all-modules-quasi-coherent\endcsname{0D1W}
\expandafter\def\csname stacksTag@etale-cohomology-lemma-category-is-colimit\endcsname{09YU}
\expandafter\def\csname stacksTag@etale-cohomology-lemma-category-loc-cst-is-colimit\endcsname{0GL2}
\expandafter\def\csname stacksTag@etale-cohomology-lemma-cech-complex\endcsname{03OZ}
\expandafter\def\csname stacksTag@etale-cohomology-lemma-characterize-locally-constant-module\endcsname{0GKC}
\expandafter\def\csname stacksTag@etale-cohomology-lemma-check-constructible\endcsname{095Q}
\expandafter\def\csname stacksTag@etale-cohomology-lemma-cofinal-etale\endcsname{03PQ}
\expandafter\def\csname stacksTag@etale-cohomology-lemma-cohomological-descent-complex-modules\endcsname{0H0U}
\expandafter\def\csname stacksTag@etale-cohomology-lemma-cohomological-dimension-proper\endcsname{095U}
\expandafter\def\csname stacksTag@etale-cohomology-lemma-cohomology-with-support-quasi-coherent\endcsname{0A46}
\expandafter\def\csname stacksTag@etale-cohomology-lemma-compare-cohomology\endcsname{0DDH}
\expandafter\def\csname stacksTag@etale-cohomology-lemma-compare-cohomology-big-small\endcsname{03YX}
\expandafter\def\csname stacksTag@etale-cohomology-lemma-compare-fppf-etale\endcsname{0F0H}
\expandafter\def\csname stacksTag@etale-cohomology-lemma-compare-h-etale\endcsname{0F0J}
\expandafter\def\csname stacksTag@etale-cohomology-lemma-compare-injectives\endcsname{0758}
\expandafter\def\csname stacksTag@etale-cohomology-lemma-compare-ph-etale\endcsname{0F0I}
\expandafter\def\csname stacksTag@etale-cohomology-lemma-compare-structure-sheaves\endcsname{07AJ}
\expandafter\def\csname stacksTag@etale-cohomology-lemma-connected-ctf-locally-constant\endcsname{09BI}
\expandafter\def\csname stacksTag@etale-cohomology-lemma-connected-locally-constant\endcsname{09BF}
\expandafter\def\csname stacksTag@etale-cohomology-lemma-constructible-abelian\endcsname{03RZ}
\expandafter\def\csname stacksTag@etale-cohomology-lemma-constructible-quasi-compact-quasi-separated\endcsname{095E}
\expandafter\def\csname stacksTag@etale-cohomology-lemma-decompose-quasi-finite-morphism\endcsname{095K}
\expandafter\def\csname stacksTag@etale-cohomology-lemma-descent-sheaf-fppf-etale\endcsname{0DEU}
\expandafter\def\csname stacksTag@etale-cohomology-lemma-describe-etale-local-ring\endcsname{04HX}
\expandafter\def\csname stacksTag@etale-cohomology-lemma-describe-pullback\endcsname{09XN}
\expandafter\def\csname stacksTag@etale-cohomology-lemma-describe-pullback-pi-fppf\endcsname{0DDL}
\expandafter\def\csname stacksTag@etale-cohomology-lemma-describe-pullback-pi-ph\endcsname{0DDW}
\expandafter\def\csname stacksTag@etale-cohomology-lemma-direct-sum-bounded-below-pushforward\endcsname{0GIW}
\expandafter\def\csname stacksTag@etale-cohomology-lemma-directed-colimit-cohomology\endcsname{03Q6}
\expandafter\def\csname stacksTag@etale-cohomology-lemma-etale-topos-independent-partial-universe\endcsname{0958}
\expandafter\def\csname stacksTag@etale-cohomology-lemma-ext-modules-hom\endcsname{0DVE}
\expandafter\def\csname stacksTag@etale-cohomology-lemma-faithful\endcsname{04LW}
\expandafter\def\csname stacksTag@etale-cohomology-lemma-finite-dim-group-cohomology\endcsname{0DVF}
\expandafter\def\csname stacksTag@etale-cohomology-lemma-finite-pushforward-constructible\endcsname{095R}
\expandafter\def\csname stacksTag@etale-cohomology-lemma-glue-etale-sheaf-fppf-cover\endcsname{0GF2}
\expandafter\def\csname stacksTag@etale-cohomology-lemma-glue-etale-sheaf-integral-surjective\endcsname{0GEZ}
\expandafter\def\csname stacksTag@etale-cohomology-lemma-glue-etale-sheaf-proper-surjective\endcsname{0GF0}
\expandafter\def\csname stacksTag@etale-cohomology-lemma-h0-proper-over-henselian-pair\endcsname{0A0C}
\expandafter\def\csname stacksTag@etale-cohomology-lemma-jshriek-constructible\endcsname{03S8}
\expandafter\def\csname stacksTag@etale-cohomology-lemma-jshriek-open\endcsname{0F70}
\expandafter\def\csname stacksTag@etale-cohomology-lemma-morphism-locally-ringed\endcsname{04I5}
\expandafter\def\csname stacksTag@etale-cohomology-lemma-points-fppf\endcsname{06VX}
\expandafter\def\csname stacksTag@etale-cohomology-lemma-points-small-etale-site\endcsname{04HU}
\expandafter\def\csname stacksTag@etale-cohomology-lemma-projection-formula-proper-mod-n\endcsname{0F0G}
\expandafter\def\csname stacksTag@etale-cohomology-lemma-proper-base-change\endcsname{0DDE}
\expandafter\def\csname stacksTag@etale-cohomology-lemma-proper-base-change-f-star\endcsname{0A3U}
\expandafter\def\csname stacksTag@etale-cohomology-lemma-proper-base-change-mod-n\endcsname{0F0C}
\expandafter\def\csname stacksTag@etale-cohomology-lemma-proper-base-change-stalk\endcsname{0DDF}
\expandafter\def\csname stacksTag@etale-cohomology-lemma-proper-mod-n-direct-sums\endcsname{0F0D}
\expandafter\def\csname stacksTag@etale-cohomology-lemma-proper-pushforward-stalk\endcsname{0A3T}
\expandafter\def\csname stacksTag@etale-cohomology-lemma-pullback-filtered-modules\endcsname{0GJ0}
\expandafter\def\csname stacksTag@etale-cohomology-lemma-pullback-on-h2-curve\endcsname{0AMB}
\expandafter\def\csname stacksTag@etale-cohomology-lemma-relative-colimit-general\endcsname{0EYM}
\expandafter\def\csname stacksTag@etale-cohomology-lemma-review-quasi-coherent\endcsname{0DEW}
\expandafter\def\csname stacksTag@etale-cohomology-lemma-sections-upstairs\endcsname{0A3H}
\expandafter\def\csname stacksTag@etale-cohomology-lemma-ses-associated-to-open\endcsname{095L}
\expandafter\def\csname stacksTag@etale-cohomology-lemma-shriek-into-star-separated-etale\endcsname{0F4L}
\expandafter\def\csname stacksTag@etale-cohomology-lemma-sp-isom-proper-loc-cst-torsion\endcsname{0GKD}
\expandafter\def\csname stacksTag@etale-cohomology-lemma-sp-isom-proper-torsion-loc-ac\endcsname{0GJW}
\expandafter\def\csname stacksTag@etale-cohomology-lemma-stalk-exact\endcsname{03PT}
\expandafter\def\csname stacksTag@etale-cohomology-lemma-stalk-pullback\endcsname{03Q1}
\expandafter\def\csname stacksTag@etale-cohomology-lemma-stalk-pushforward-closed-immersion\endcsname{04FX}
\expandafter\def\csname stacksTag@etale-cohomology-lemma-support-section-closed\endcsname{04FT}
\expandafter\def\csname stacksTag@etale-cohomology-lemma-supported-in-closed-points\endcsname{0F1F}
\expandafter\def\csname stacksTag@etale-cohomology-lemma-tensor-c\endcsname{0961}
\expandafter\def\csname stacksTag@etale-cohomology-lemma-tensor-constructible\endcsname{0GKB}
\expandafter\def\csname stacksTag@etale-cohomology-lemma-torsion-direct-image\endcsname{0DDD}
\expandafter\def\csname stacksTag@etale-cohomology-lemma-what-integral\endcsname{04C2}
\expandafter\def\csname stacksTag@etale-cohomology-lemma-when-ctf\endcsname{03TT}
\expandafter\def\csname stacksTag@etale-cohomology-lemma-zero-over-image\endcsname{04FR}
\expandafter\def\csname stacksTag@etale-cohomology-proposition-closed-immersion-pushforward\endcsname{04CA}
\expandafter\def\csname stacksTag@etale-cohomology-proposition-constructible-over-noetherian\endcsname{09BH}
\expandafter\def\csname stacksTag@etale-cohomology-proposition-describe-jshriek\endcsname{03S5}
\expandafter\def\csname stacksTag@etale-cohomology-proposition-finite-higher-direct-image-zero\endcsname{03QP}
\expandafter\def\csname stacksTag@etale-cohomology-proposition-integral-universally-injective-pushforward\endcsname{04FZ}
\expandafter\def\csname stacksTag@etale-cohomology-proposition-quasi-coherent-sheaf-fpqc\endcsname{03OG}
\expandafter\def\csname stacksTag@etale-cohomology-remark-affine-inside-equivalence\endcsname{05YX}
\expandafter\def\csname stacksTag@etale-cohomology-remark-constant-locally-constant-maps\endcsname{03P5}
\expandafter\def\csname stacksTag@etale-cohomology-remark-stalk-pullback\endcsname{04JN}
\expandafter\def\csname stacksTag@etale-cohomology-remarks-enough-points\endcsname{040S}
\expandafter\def\csname stacksTag@etale-cohomology-section-Dc\endcsname{095V}
\expandafter\def\csname stacksTag@etale-cohomology-section-artin-schreier\endcsname{0A3J}
\expandafter\def\csname stacksTag@etale-cohomology-section-base-change-preliminaries\endcsname{0EZQ}
\expandafter\def\csname stacksTag@etale-cohomology-section-cohomology-support\endcsname{09XP}
\expandafter\def\csname stacksTag@etale-cohomology-section-compare\endcsname{0757}
\expandafter\def\csname stacksTag@etale-cohomology-section-compare-topologies\endcsname{09XL}
\expandafter\def\csname stacksTag@etale-cohomology-section-computation\endcsname{03N8}
\expandafter\def\csname stacksTag@etale-cohomology-section-continuous-group-cohomology\endcsname{0DVG}
\expandafter\def\csname stacksTag@etale-cohomology-section-extension-by-zero\endcsname{03S2}
\expandafter\def\csname stacksTag@etale-cohomology-section-fppf-etale\endcsname{0DDK}
\expandafter\def\csname stacksTag@etale-cohomology-section-functoriality\endcsname{04I0}
\expandafter\def\csname stacksTag@etale-cohomology-section-galois-action-stalks\endcsname{03QW}
\expandafter\def\csname stacksTag@etale-cohomology-section-glueing-etale\endcsname{0GEX}
\expandafter\def\csname stacksTag@etale-cohomology-section-introduction\endcsname{03N2}
\expandafter\def\csname stacksTag@etale-cohomology-section-ph-etale\endcsname{0DDV}
\expandafter\def\csname stacksTag@etale-cohomology-section-proper-base-change\endcsname{095S}
\expandafter\def\csname stacksTag@etale-cohomology-section-stalks\endcsname{03PN}
\expandafter\def\csname stacksTag@etale-cohomology-section-stalks-structure-sheaf\endcsname{04HW}
\expandafter\def\csname stacksTag@etale-cohomology-section-trace-method\endcsname{03SH}
\expandafter\def\csname stacksTag@etale-cohomology-section-vanishing-torsion\endcsname{03SB}
\expandafter\def\csname stacksTag@etale-cohomology-theorem-colimit\endcsname{09YQ}
\expandafter\def\csname stacksTag@etale-cohomology-theorem-equivalence-sheaves-point\endcsname{03QT}
\expandafter\def\csname stacksTag@etale-cohomology-theorem-exactness-stalks\endcsname{03PU}
\expandafter\def\csname stacksTag@etale-cohomology-theorem-fully-faithful\endcsname{04I7}
\expandafter\def\csname stacksTag@etale-cohomology-theorem-higher-direct-images\endcsname{03Q9}
\expandafter\def\csname stacksTag@etale-cohomology-theorem-proper-base-change\endcsname{095T}
\expandafter\def\csname stacksTag@etale-cohomology-theorem-quasi-coherent\endcsname{03OJ}
\expandafter\def\csname stacksTag@etale-cohomology-theorem-topological-invariance\endcsname{04DZ}
\expandafter\def\csname stacksTag@etale-cohomology-theorem-vanishing-affine-curves\endcsname{03SC}
\expandafter\def\csname stacksTag@etale-cohomology-theorem-vanishing-affine-curves-coefficients\endcsname{0GJI}
\expandafter\def\csname stacksTag@pione-definition-G-set-continuous\endcsname{04JJ}
\expandafter\def\csname stacksTag@pione-definition-fundamental-group\endcsname{0BNC}
\expandafter\def\csname stacksTag@pione-definition-galois-category\endcsname{0BMY}
\expandafter\def\csname stacksTag@pione-lemma-inertial-invariants-unramified\endcsname{0BSU}
\expandafter\def\csname stacksTag@pione-section-finite-etale\endcsname{0BL6}
\expandafter\def\csname stacksTag@pione-section-group-actions-integral\endcsname{0BSN}
\expandafter\def\csname stacksTag@pione-section-introduction\endcsname{0BQ7}
\expandafter\def\csname stacksTag@spaces-more-groupoids-definition-split-at-point\endcsname{04RK}
\expandafter\def\csname stacksTag@spaces-more-groupoids-lemma-composition-is-addition\endcsname{0CKF}
\expandafter\def\csname stacksTag@spaces-more-groupoids-lemma-group-scheme-over-field-separated\endcsname{08BH}
\expandafter\def\csname stacksTag@spaces-more-groupoids-lemma-group-space-scheme-locally-finite-type-over-k\endcsname{0B8F}
\expandafter\def\csname stacksTag@spaces-more-groupoids-lemma-groupoid-on-field-dimension-equal-stabilizer\endcsname{06FF}
\expandafter\def\csname stacksTag@spaces-more-groupoids-lemma-no-nonconstant-morphism-from-P1-to-group\endcsname{0AEN}
\expandafter\def\csname stacksTag@spaces-more-groupoids-lemma-property-G-invariant\endcsname{06R4}
\expandafter\def\csname stacksTag@spaces-more-groupoids-lemma-quasi-splitting-affine-scheme\endcsname{04S0}
\expandafter\def\csname stacksTag@spaces-more-groupoids-lemma-restrict-preserves-type\endcsname{04RP}
\expandafter\def\csname stacksTag@spaces-more-groupoids-lemma-splitting-affine-scheme\endcsname{0DTE}
\expandafter\def\csname stacksTag@spaces-more-groupoids-lemma-strong-splitting-affine-scheme\endcsname{04RZ}
\expandafter\def\csname stacksTag@spaces-more-groupoids-proposition-finite-algebraic-space\endcsname{04QH}
\expandafter\def\csname stacksTag@spaces-more-groupoids-proposition-group-space-scheme-over-field\endcsname{0B8G}
\expandafter\def\csname stacksTag@spaces-more-groupoids-section-etale-localize\endcsname{04RJ}
\expandafter\def\csname stacksTag@spaces-more-groupoids-section-finite\endcsname{04PB}
\expandafter\def\csname stacksTag@spaces-more-groupoids-section-groupoid-sections\endcsname{0CKB}
\expandafter\def\csname stacksTag@varieties-definition-absolute-frobenius\endcsname{03SM}
\expandafter\def\csname stacksTag@varieties-definition-algebraic-scheme\endcsname{06LG}
\expandafter\def\csname stacksTag@varieties-definition-curve\endcsname{0A23}
\expandafter\def\csname stacksTag@varieties-definition-degree\endcsname{0BEW}
\expandafter\def\csname stacksTag@varieties-definition-degree-invertible-sheaf\endcsname{0AYR}
\expandafter\def\csname stacksTag@varieties-definition-delta-invariant\endcsname{0C1T}
\expandafter\def\csname stacksTag@varieties-definition-delta-invariant-algebra\endcsname{0C3T}
\expandafter\def\csname stacksTag@varieties-definition-dual-numbers\endcsname{0B29}
\expandafter\def\csname stacksTag@varieties-definition-embed-dim\endcsname{0C1Q}
\expandafter\def\csname stacksTag@varieties-definition-embedding-dimension\endcsname{0C2H}
\expandafter\def\csname stacksTag@varieties-definition-euler-characteristic\endcsname{0BEJ}
\expandafter\def\csname stacksTag@varieties-definition-geometrically-connected\endcsname{0362}
\expandafter\def\csname stacksTag@varieties-definition-geometrically-integral\endcsname{020H}
\expandafter\def\csname stacksTag@varieties-definition-geometrically-irreducible\endcsname{0365}
\expandafter\def\csname stacksTag@varieties-definition-geometrically-normal\endcsname{038M}
\expandafter\def\csname stacksTag@varieties-definition-geometrically-reduced\endcsname{035V}
\expandafter\def\csname stacksTag@varieties-definition-geometrically-regular\endcsname{038T}
\expandafter\def\csname stacksTag@varieties-definition-hilbert-polynomial\endcsname{08AD}
\expandafter\def\csname stacksTag@varieties-definition-intersection-number\endcsname{0BEP}
\expandafter\def\csname stacksTag@varieties-definition-regularity\endcsname{08A3}
\expandafter\def\csname stacksTag@varieties-definition-relative-frobenius\endcsname{0CC9}
\expandafter\def\csname stacksTag@varieties-definition-tangent-space\endcsname{0B2C}
\expandafter\def\csname stacksTag@varieties-definition-variety\endcsname{020D}
\expandafter\def\csname stacksTag@varieties-definition-variety-type\endcsname{04L1}
\expandafter\def\csname stacksTag@varieties-definition-wedge\endcsname{0C41}
\expandafter\def\csname stacksTag@varieties-equation-tangent-space-fibre\endcsname{0BEA}
\expandafter\def\csname stacksTag@varieties-lemma-CM-base-change\endcsname{045P}
\expandafter\def\csname stacksTag@varieties-lemma-Galois-action-quasi-compact-open\endcsname{04KU}
\expandafter\def\csname stacksTag@varieties-lemma-affine-geometrically-connected\endcsname{0386}
\expandafter\def\csname stacksTag@varieties-lemma-affine-space-over-field\endcsname{055T}
\expandafter\def\csname stacksTag@varieties-lemma-algebraic-scheme-dim-0\endcsname{06LH}
\expandafter\def\csname stacksTag@varieties-lemma-ample-after-field-extension\endcsname{0BDC}
\expandafter\def\csname stacksTag@varieties-lemma-ample-curve\endcsname{0B5X}
\expandafter\def\csname stacksTag@varieties-lemma-ample-positive\endcsname{0BEV}
\expandafter\def\csname stacksTag@varieties-lemma-ampleness-in-terms-of-degrees-components\endcsname{0B5Y}
\expandafter\def\csname stacksTag@varieties-lemma-automorphism\endcsname{0G05}
\expandafter\def\csname stacksTag@varieties-lemma-baby-stein\endcsname{0FD1}
\expandafter\def\csname stacksTag@varieties-lemma-base-change-complete-local-ring\endcsname{0C54}
\expandafter\def\csname stacksTag@varieties-lemma-bertini\endcsname{0FD6}
\expandafter\def\csname stacksTag@varieties-lemma-bijection-connected-components\endcsname{0385}
\expandafter\def\csname stacksTag@varieties-lemma-bijection-irreducible-components\endcsname{038F}
\expandafter\def\csname stacksTag@varieties-lemma-bound-quotients-free\endcsname{08AG}
\expandafter\def\csname stacksTag@varieties-lemma-change-fields-pic\endcsname{0CC5}
\expandafter\def\csname stacksTag@varieties-lemma-char-p-differentials-free-smooth\endcsname{04QP}
\expandafter\def\csname stacksTag@varieties-lemma-char-zero-differentials-free-smooth\endcsname{04QN}
\expandafter\def\csname stacksTag@varieties-lemma-characterize-geometrically-connected\endcsname{0387}
\expandafter\def\csname stacksTag@varieties-lemma-characterize-geometrically-disconnected\endcsname{0389}
\expandafter\def\csname stacksTag@varieties-lemma-check-invertible-sheaf-trivial\endcsname{0B40}
\expandafter\def\csname stacksTag@varieties-lemma-chi-tensor-finite\endcsname{0AYT}
\expandafter\def\csname stacksTag@varieties-lemma-closed-fixed-by-Galois\endcsname{038B}
\expandafter\def\csname stacksTag@varieties-lemma-closure-image-product-map\endcsname{04Q0}
\expandafter\def\csname stacksTag@varieties-lemma-closure-of-product\endcsname{047B}
\expandafter\def\csname stacksTag@varieties-lemma-complement-codim-1-closed-points\endcsname{09NA}
\expandafter\def\csname stacksTag@varieties-lemma-complete-local-ring-structure-as-algebra\endcsname{0C52}
\expandafter\def\csname stacksTag@varieties-lemma-connectedness-ample-divisor\endcsname{0FD9}
\expandafter\def\csname stacksTag@varieties-lemma-degree-additive\endcsname{0AYS}
\expandafter\def\csname stacksTag@varieties-lemma-degree-base-change\endcsname{0B59}
\expandafter\def\csname stacksTag@varieties-lemma-degree-birational-pullback\endcsname{0AYU}
\expandafter\def\csname stacksTag@varieties-lemma-degree-effective-Cartier-divisor\endcsname{0AYY}
\expandafter\def\csname stacksTag@varieties-lemma-degree-finite-morphism-in-terms-degrees\endcsname{0BEX}
\expandafter\def\csname stacksTag@varieties-lemma-degree-in-terms-of-components\endcsname{0AYW}
\expandafter\def\csname stacksTag@varieties-lemma-degree-on-proper-curve\endcsname{0AYV}
\expandafter\def\csname stacksTag@varieties-lemma-degree-pullback-map-proper-curves\endcsname{0AYZ}
\expandafter\def\csname stacksTag@varieties-lemma-degree-tensor-product\endcsname{0AYX}
\expandafter\def\csname stacksTag@varieties-lemma-delta-invariant-and-change-of-fields\endcsname{0C3X}
\expandafter\def\csname stacksTag@varieties-lemma-delta-invariant-and-change-of-fields-better\endcsname{0C3Y}
\expandafter\def\csname stacksTag@varieties-lemma-delta-invariant-is-zero\endcsname{0C3U}
\expandafter\def\csname stacksTag@varieties-lemma-delta-number-branches-inequality\endcsname{0C43}
\expandafter\def\csname stacksTag@varieties-lemma-delta-number-branches-inequality-sh\endcsname{0C42}
\expandafter\def\csname stacksTag@varieties-lemma-delta-same-after-completion\endcsname{0C3W}
\expandafter\def\csname stacksTag@varieties-lemma-dim-1-projective-completion\endcsname{0BXV}
\expandafter\def\csname stacksTag@varieties-lemma-dim-1-projective-completion-after-insep\endcsname{0GK5}
\expandafter\def\csname stacksTag@varieties-lemma-dim-1-proper-projective\endcsname{0A26}
\expandafter\def\csname stacksTag@varieties-lemma-dim-1-quasi-projective\endcsname{0A25}
\expandafter\def\csname stacksTag@varieties-lemma-dimension-fibre-in-higher-codimension\endcsname{0B2K}
\expandafter\def\csname stacksTag@varieties-lemma-dimension-fibres-locally-algebraic\endcsname{0B2L}
\expandafter\def\csname stacksTag@varieties-lemma-dimension-locally-algebraic\endcsname{0A21}
\expandafter\def\csname stacksTag@varieties-lemma-dimension-product-locally-algebraic\endcsname{0B2M}
\expandafter\def\csname stacksTag@varieties-lemma-divisible\endcsname{0C6P}
\expandafter\def\csname stacksTag@varieties-lemma-dominate-valuation-ring-dimension-fibres\endcsname{0B2J}
\expandafter\def\csname stacksTag@varieties-lemma-equation-codim-1-in-projective-space\endcsname{0BXU}
\expandafter\def\csname stacksTag@varieties-lemma-euler-characteristic-additive\endcsname{08AA}
\expandafter\def\csname stacksTag@varieties-lemma-exact-sequence-induction\endcsname{08A0}
\expandafter\def\csname stacksTag@varieties-lemma-finite-extension-geometrically-irreducible-components\endcsname{054R}
\expandafter\def\csname stacksTag@varieties-lemma-finite-extension-geometrically-normal\endcsname{0BXT}
\expandafter\def\csname stacksTag@varieties-lemma-finite-extension-geometrically-reduced\endcsname{04KT}
\expandafter\def\csname stacksTag@varieties-lemma-finite-in-codim-1\endcsname{0AB7}
\expandafter\def\csname stacksTag@varieties-lemma-flat-under-smooth\endcsname{05AX}
\expandafter\def\csname stacksTag@varieties-lemma-general-degree-g-line-bundle\endcsname{0B8Z}
\expandafter\def\csname stacksTag@varieties-lemma-generate-over-complement\endcsname{0FD5}
\expandafter\def\csname stacksTag@varieties-lemma-generic-points-geometrically-reduced\endcsname{04KS}
\expandafter\def\csname stacksTag@varieties-lemma-geometric-branches-and-change-of-fields\endcsname{0C40}
\expandafter\def\csname stacksTag@varieties-lemma-geometric-structure-unramified\endcsname{0CBJ}
\expandafter\def\csname stacksTag@varieties-lemma-geometrically-connected-check-after-extension\endcsname{054N}
\expandafter\def\csname stacksTag@varieties-lemma-geometrically-connected-criterion\endcsname{056R}
\expandafter\def\csname stacksTag@varieties-lemma-geometrically-connected-if-connected-and-point\endcsname{04KV}
\expandafter\def\csname stacksTag@varieties-lemma-geometrically-integral\endcsname{038K}
\expandafter\def\csname stacksTag@varieties-lemma-geometrically-irreducible-check-after-extension\endcsname{054P}
\expandafter\def\csname stacksTag@varieties-lemma-geometrically-irreducible-function-field\endcsname{054Q}
\expandafter\def\csname stacksTag@varieties-lemma-geometrically-normal-in-codim-1\endcsname{0C3P}
\expandafter\def\csname stacksTag@varieties-lemma-geometrically-normal-stein\endcsname{0FD3}
\expandafter\def\csname stacksTag@varieties-lemma-geometrically-normal-upstairs\endcsname{038P}
\expandafter\def\csname stacksTag@varieties-lemma-geometrically-reduced\endcsname{035X}
\expandafter\def\csname stacksTag@varieties-lemma-geometrically-reduced-any-base-change\endcsname{035Z}
\expandafter\def\csname stacksTag@varieties-lemma-geometrically-reduced-at-point\endcsname{035W}
\expandafter\def\csname stacksTag@varieties-lemma-geometrically-reduced-dense-smooth-open\endcsname{056V}
\expandafter\def\csname stacksTag@varieties-lemma-geometrically-reduced-upstairs\endcsname{0384}
\expandafter\def\csname stacksTag@varieties-lemma-geometrically-regular\endcsname{038V}
\expandafter\def\csname stacksTag@varieties-lemma-geometrically-regular-smooth\endcsname{038X}
\expandafter\def\csname stacksTag@varieties-lemma-geometrically-regular-upstairs\endcsname{038W}
\expandafter\def\csname stacksTag@varieties-lemma-geometrically-unibranch-and-change-of-fields\endcsname{0C55}
\expandafter\def\csname stacksTag@varieties-lemma-globally-generated-base-change\endcsname{0B57}
\expandafter\def\csname stacksTag@varieties-lemma-image-connected-component\endcsname{04PZ}
\expandafter\def\csname stacksTag@varieties-lemma-image-irreducible\endcsname{04KX}
\expandafter\def\csname stacksTag@varieties-lemma-immersion-into-affine\endcsname{0CBL}
\expandafter\def\csname stacksTag@varieties-lemma-injective-tangent-spaces-unramified\endcsname{0B2G}
\expandafter\def\csname stacksTag@varieties-lemma-inseparable-deg-p-smooth\endcsname{0CCF}
\expandafter\def\csname stacksTag@varieties-lemma-integral-closure-ground-field\endcsname{04MI}
\expandafter\def\csname stacksTag@varieties-lemma-intersection-in-projective-space\endcsname{0B2R}
\expandafter\def\csname stacksTag@varieties-lemma-intersection-number-additive\endcsname{0BER}
\expandafter\def\csname stacksTag@varieties-lemma-intersection-number-and-pullback\endcsname{0BET}
\expandafter\def\csname stacksTag@varieties-lemma-intersection-numbers-and-degrees-on-curves\endcsname{0BEY}
\expandafter\def\csname stacksTag@varieties-lemma-irreducible-components-geometrically-irreducible\endcsname{054S}
\expandafter\def\csname stacksTag@varieties-lemma-kunneth\endcsname{0BED}
\expandafter\def\csname stacksTag@varieties-lemma-locally-Noetherian-base-change\endcsname{038R}
\expandafter\def\csname stacksTag@varieties-lemma-locally-finite-type-Jacobson\endcsname{0478}
\expandafter\def\csname stacksTag@varieties-lemma-m-regular-globally-generated\endcsname{08A8}
\expandafter\def\csname stacksTag@varieties-lemma-m-regular-up\endcsname{08A6}
\expandafter\def\csname stacksTag@varieties-lemma-make-Jacobson\endcsname{0479}
\expandafter\def\csname stacksTag@varieties-lemma-map-tangent-spaces\endcsname{0B2F}
\expandafter\def\csname stacksTag@varieties-lemma-modification-normal-iso-over-codimension-1\endcsname{0BFP}
\expandafter\def\csname stacksTag@varieties-lemma-no-sections-dual-nef\endcsname{0E22}
\expandafter\def\csname stacksTag@varieties-lemma-normalization-locally-algebraic\endcsname{0BXR}
\expandafter\def\csname stacksTag@varieties-lemma-normalization-same-after-completion\endcsname{0C3V}
\expandafter\def\csname stacksTag@varieties-lemma-numerical-intersection-effective-Cartier-divisor\endcsname{0BEU}
\expandafter\def\csname stacksTag@varieties-lemma-numerical-polynomial-from-euler\endcsname{0BEM}
\expandafter\def\csname stacksTag@varieties-lemma-numerical-polynomial-leading-term\endcsname{0BEN}
\expandafter\def\csname stacksTag@varieties-lemma-orbit-irreducible-components\endcsname{04KY}
\expandafter\def\csname stacksTag@varieties-lemma-perfect-reduced\endcsname{020I}
\expandafter\def\csname stacksTag@varieties-lemma-pre-delta-invariant\endcsname{0C3S}
\expandafter\def\csname stacksTag@varieties-lemma-prepare-delta-invariant\endcsname{0C1R}
\expandafter\def\csname stacksTag@varieties-lemma-product-varieties\endcsname{05P3}
\expandafter\def\csname stacksTag@varieties-lemma-proper-geometrically-reduced-global-sections\endcsname{0BUG}
\expandafter\def\csname stacksTag@varieties-lemma-quasi-finite-in-codim-1\endcsname{0AB6}
\expandafter\def\csname stacksTag@varieties-lemma-rational-equivalence-for-Pic\endcsname{0BEH}
\expandafter\def\csname stacksTag@varieties-lemma-reduced-dim-1-projective-completion\endcsname{0BXW}
\expandafter\def\csname stacksTag@varieties-lemma-regular-functions-proper-variety\endcsname{04L2}
\expandafter\def\csname stacksTag@varieties-lemma-regular-point-on-curve\endcsname{0B8Y}
\expandafter\def\csname stacksTag@varieties-lemma-relative-frobenius\endcsname{0CCB}
\expandafter\def\csname stacksTag@varieties-lemma-relative-frobenius-endomorphism-identity\endcsname{0CCA}
\expandafter\def\csname stacksTag@varieties-lemma-relative-frobenius-finite\endcsname{0CCD}
\expandafter\def\csname stacksTag@varieties-lemma-relative-frobenius-omega\endcsname{0CCC}
\expandafter\def\csname stacksTag@varieties-lemma-scheme-theoretic-image-product-map\endcsname{04Q1}
\expandafter\def\csname stacksTag@varieties-lemma-separably-closed-field-connected-components\endcsname{0363}
\expandafter\def\csname stacksTag@varieties-lemma-separably-closed-field-irreducible-components\endcsname{038H}
\expandafter\def\csname stacksTag@varieties-lemma-separably-closed-irreducible\endcsname{020J}
\expandafter\def\csname stacksTag@varieties-lemma-smooth-geometrically-normal\endcsname{056T}
\expandafter\def\csname stacksTag@varieties-lemma-smooth-regular\endcsname{056S}
\expandafter\def\csname stacksTag@varieties-lemma-smooth-separable-closed-points-dense\endcsname{056U}
\expandafter\def\csname stacksTag@varieties-lemma-smooth-stratification\endcsname{0H3W}
\expandafter\def\csname stacksTag@varieties-lemma-surjective-pic-birational-finite\endcsname{0C0T}
\expandafter\def\csname stacksTag@varieties-lemma-tangent-space-cotangent-space\endcsname{0B2D}
\expandafter\def\csname stacksTag@varieties-lemma-tangent-space-product\endcsname{0BEB}
\expandafter\def\csname stacksTag@varieties-lemma-vanishin-h0-negative\endcsname{0FD7}
\expandafter\def\csname stacksTag@varieties-lemma-vanishing-degree-2g-and-1-line-bundle\endcsname{0B90}
\expandafter\def\csname stacksTag@varieties-lemma-variant-separate-points-tangent-vectors\endcsname{0E8T}
\expandafter\def\csname stacksTag@varieties-lemma-variety-with-smooth-rational-point\endcsname{0CDW}
\expandafter\def\csname stacksTag@varieties-lemma-very-ample-vanish-at-point\endcsname{0B58}
\expandafter\def\csname stacksTag@varieties-proposition-finite-set-of-points-of-codim-1-in-affine\endcsname{09NN}
\expandafter\def\csname stacksTag@varieties-proposition-unique-base-field\endcsname{04MK}
\expandafter\def\csname stacksTag@varieties-remark-n-fold-relative-frobenius\endcsname{0CCG}
\expandafter\def\csname stacksTag@varieties-section-CM\endcsname{045O}
\expandafter\def\csname stacksTag@varieties-section-algebraic-schemes\endcsname{06LF}
\expandafter\def\csname stacksTag@varieties-section-bertini\endcsname{0FD4}
\expandafter\def\csname stacksTag@varieties-section-curves\endcsname{0A22}
\expandafter\def\csname stacksTag@varieties-section-dimension-one\endcsname{09N7}
\expandafter\def\csname stacksTag@varieties-section-divisors-curves\endcsname{0AYQ}
\expandafter\def\csname stacksTag@varieties-section-embedding-dimension\endcsname{0C2G}
\expandafter\def\csname stacksTag@varieties-section-frobenius\endcsname{0CC6}
\expandafter\def\csname stacksTag@varieties-section-generically-finite\endcsname{0AB5}
\expandafter\def\csname stacksTag@varieties-section-kunneth\endcsname{0BEC}
\expandafter\def\csname stacksTag@varieties-section-normalization\endcsname{0BXQ}
\expandafter\def\csname stacksTag@varieties-section-normalization-one-dimensional\endcsname{0C44}
\expandafter\def\csname stacksTag@varieties-section-num\endcsname{0BEL}
\expandafter\def\csname stacksTag@varieties-section-number-of-branches\endcsname{0C3Z}
\expandafter\def\csname stacksTag@varieties-section-picard-groups\endcsname{0BEG}
\expandafter\def\csname stacksTag@varieties-section-separating-points-tangent-vectors\endcsname{0E8R}
\expandafter\def\csname stacksTag@varieties-section-smooth\endcsname{04QM}
\expandafter\def\csname stacksTag@varieties-section-varieties\endcsname{020C}
\expandafter\def\csname stacksTag@varieties-section-varieties-rational-maps\endcsname{0BXM}
\expandafter\def\csname stacksTag@varieties-situation-family-divisors\endcsname{0G47}
\expandafter\def\csname stacksTag@varieties-subsection-hilbert\endcsname{08A9}
\expandafter\def\csname stacksTag@varieties-subsection-regularity\endcsname{08A2}
\expandafter\def\csname stacksTag@varieties-theorem-varieties-rational-maps\endcsname{0BXN}
\expandafter\def\csname stacksManual@foo-lemma-bar\endcsname{\href{https://github.com/stacks/stacks-project/blob/a04446e57ec1fbc252a871afcec7752fb2807b14/coding.tex}{\texttt{源码示例}}}
\DeclareRobustCommand{\StacksResolvedRef}[1]{%
  \ifcsname r@#1\endcsname
    \StacksOriginalRef{#1}%
  \else
    \ifcsname stacksTag@#1\endcsname
      \href{https://stacks.math.columbia.edu/tag/\csname stacksTag@#1\endcsname}{\texttt{\csname stacksTag@#1\endcsname}}%
    \else
      \ifcsname stacksManual@#1\endcsname
        \csname stacksManual@#1\endcsname%
      \else
        \StacksOriginalRef{#1}%
      \fi
    \fi
  \fi
}
\AtBeginDocument{%
  \let\StacksOriginalRef\ref
  \let\ref\StacksResolvedRef
}


% Theorem environments.
%
\theoremstyle{plain}
\newtheorem{theorem}[subsection]{定理}
\newtheorem{proposition}[subsection]{命题}
\newtheorem{lemma}[subsection]{引理}

\theoremstyle{definition}
\newtheorem{definition}[subsection]{定义}
\newtheorem{example}[subsection]{例}
\newtheorem{exercise}[subsection]{习题}
\newtheorem{situation}[subsection]{情形}

\theoremstyle{remark}
\newtheorem{remark}[subsection]{注}
\newtheorem{remarks}[subsection]{诸注}

\numberwithin{equation}{subsection}

% Macros
%
\def\lim{\mathop{\mathrm{lim}}\nolimits}
\def\colim{\mathop{\mathrm{colim}}\nolimits}
\def\Spec{\mathop{\mathrm{Spec}}}
\def\Hom{\mathop{\mathrm{Hom}}\nolimits}
\def\Ext{\mathop{\mathrm{Ext}}\nolimits}
\def\SheafHom{\mathop{\mathcal{H}\!\mathit{om}}\nolimits}
\def\SheafExt{\mathop{\mathcal{E}\!\mathit{xt}}\nolimits}
\def\Sch{\mathit{Sch}}
\def\Mor{\mathop{\mathrm{Mor}}\nolimits}
\def\Ob{\mathop{\mathrm{Ob}}\nolimits}
\def\Sh{\mathop{\mathit{Sh}}\nolimits}
\def\NL{\mathop{N\!L}\nolimits}
\def\CH{\mathop{\mathrm{CH}}\nolimits}
\def\proetale{{pro\text{-}\acute{e}tale}}
\def\etale{{\acute{e}tale}}
\def\QCoh{\mathit{QCoh}}
\def\Ker{\mathop{\mathrm{Ker}}}
\def\Im{\mathop{\mathrm{Im}}}
\def\Coker{\mathop{\mathrm{Coker}}}
\def\Coim{\mathop{\mathrm{Coim}}}

% Boxtimes
%
\DeclareMathSymbol{\boxtimes}{\mathbin}{AMSa}{"02}

%
% Macros for moduli stacks/spaces
%
\def\QCohstack{\mathcal{QC}\!\mathit{oh}}
\def\Cohstack{\mathcal{C}\!\mathit{oh}}
\def\Spacesstack{\mathcal{S}\!\mathit{paces}}
\def\Quotfunctor{\mathrm{Quot}}
\def\Hilbfunctor{\mathrm{Hilb}}
\def\Curvesstack{\mathcal{C}\!\mathit{urves}}
\def\Polarizedstack{\mathcal{P}\!\mathit{olarized}}
\def\Complexesstack{\mathcal{C}\!\mathit{omplexes}}
% \Pic is the operator that assigns to X its picard group, usage \Pic(X)
% \Picardstack_{X/B} denotes the Picard stack of X over B
% \Picardfunctor_{X/B} denotes the Picard functor of X over B
\def\Pic{\mathop{\mathrm{Pic}}\nolimits}
\def\Picardstack{\mathcal{P}\!\mathit{ic}}
\def\Picardfunctor{\mathrm{Pic}}
\def\Deformationcategory{\mathcal{D}\!\mathit{ef}}
